\subsection{Herramientas y Tec\-no\-lo\-gí\-as}
\label{sec:tools_technologies}

Se utilizará Python 3.11 como lenguaje de programación base del desarrollo, selección justificada por la disponibilidad de un ecosistema extenso de bibliotecas orientadas al aprendizaje automático y aprendizaje profundo. Conforme a Rana y Bhushan \cite{Rana2022}, Python se ha posicionado como el lenguaje predominante en aplicaciones de inteligencia artificial centradas en análisis de imágenes médicas. Además, se empleará Google Colab como entorno computacional en línea, proporcionando acceso a aceleradores GPU y facilitando la replicabilidad de experimentos sin requerer infraestructura local de elevado desempeño. Según Puttagunta y Ravi \cite{Puttagunta2021}, la configuración adecuada del ambiente de hardware y software es fundamental para asegurar la reproducibilidad de investigaciones computacionales.

Se empleará TensorFlow 2.18.0 y Keras 3.8.0 para la implementación de arquitecturas basadas en convoluciones tales como ResNet50 y EfficientNet-B0, en tanto que los modelos fundamentados en Transformers (ViT-B/16, Swin Transformer y CoAtNet) serán programados utilizando PyTorch 2.6 y timm 1.0. Beirović et al. \cite{Beirovic2025} presentan un análisis comparativo de estos frameworks para categorización de imágenes médicas. El acondicionamiento previo de las imágenes se realizará con OpenCV 4.10, NumPy 2.2 para cálculos matriciales, scikit-learn 1.6 para métricas de desempeño y particionamiento de datos, y Matplotlib 3.10 para visualización de resultados. Laçi et al. \cite{Laci2025} señalan que TensorFlow, Keras y PyTorch constituyen los frameworks más ampliamente adoptados en aprendizaje profundo aplicado a imágenes médicas.
